% =========================================================================
% PLANTILLA MAESTRA: TRANSCRIPCIÓN DE DIAPOSITIVAS Y TEXTO ULTRA-DENSO
% =========================================================================
% Optimizada para vaciar presentaciones de PowerPoint (PPT/PDF de 30-80 filminas),
% guías infladas o transcripciones completas sin resumir, maximizando la
% superficie útil de papel A4 y reduciendo drásticamente el número de páginas.
%
% Principios de Diseño Editorial:
% 1. Formato a 2 Columnas Continuas (twocolumn): Longitud de línea ergonómica
%    (55-65 caracteres) perfecta para viñetas breves y oraciones cortas de PPT.
% 2. Geometría Extrema (Márgenes de 1.0 cm): Cero desperdicio de papel (92% de
%    superficie útil aprovechada).
% 3. Cero Desperdicio Inicial: Sin portada separada, sin índice y sin banners
%    pesados. El texto comienza en el renglón 1 de la página 1.
% 4. Macro \filmina{N}{Título}: Marcador de diapositiva inline sin salto de
%    página, permitiendo ubicar el contenido de la presentación en lectura continua.
% 5. Tipografía Láser Monocromo 600 DPI: Bitstream Charter a 9pt con microtype
%    para empaquetado perfecto sin huérfanas ni desbordes.
% =========================================================================

\documentclass[9pt,a4paper,twocolumn]{extarticle}

% --- Codificación y Lenguaje ---
\usepackage[utf8]{inputenc}
\usepackage[T1]{fontenc}
\usepackage[spanish,es-tabla,es-noquoting]{babel}

% --- Tipografía Láser de Alto Rendimiento ---
\usepackage{charter}                         % Bitstream Charter: Excelente visibilidad con tóner
\usepackage{tgheros}                         % TeX Gyre Heros para títulos sans-serif
\usepackage{courier}                         % Courier para parámetros y comandos monoespaciados
\usepackage{microtype}                       % Justificación óptica de alta densidad

% --- Geometría Extrema ---
\usepackage{geometry}
\geometry{
    top=1.0cm,
    bottom=1.1cm,
    left=1.0cm,
    right=1.0cm,
    columnsep=6mm,                           % Separación compacta entre columnas
    headheight=0pt,                          % Cero encabezado superior
    headsep=0pt,                             % Cero separación de encabezado
    footskip=6mm                             % Pie de página ultra compacto
}

% Regla divisoria vertical sutil entre columnas
\setlength{\columnseprule}{0.3pt}

% --- Interlineado y Párrafos ---
\linespread{1.08}
\setlength{\parindent}{0pt}
\setlength{\parskip}{2pt}
\raggedbottom

% --- Paquetes Complementarios ---
\usepackage{amsmath,amssymb}
\usepackage{booktabs}
\usepackage{tabularx}
\usepackage{enumitem}
\usepackage{listings}
\usepackage{xcolor}
\usepackage{lastpage}
\usepackage{hyperref}

% --- Listas Ultra-Compactas para Viñetas de PPT ---
\setlist[itemize]{leftmargin=1.1em, itemsep=1.2pt, parsep=0pt, topsep=1.5pt, label=\raisebox{0.15ex}{\tiny$\blacktriangleright$}}
\setlist[enumerate]{leftmargin=1.3em, itemsep=1.2pt, parsep=0pt, topsep=1.5pt}

% --- Configuración de Código / Comandos ---
\lstdefinestyle{compactcode}{
    basicstyle=\ttfamily\fontsize{7.5pt}{9pt}\selectfont\color{black},
    keywordstyle=\bfseries\color{black},
    commentstyle=\itshape\color{black!70},
    stringstyle=\color{black!90},
    showstringspaces=false,
    frame=single,
    framerule=0.3pt,
    rulecolor=\color{black!60},
    backgroundcolor=\color{white},
    breaklines=true,
    tabsize=2,
    aboveskip=3pt,
    belowskip=3pt,
    xleftmargin=4pt,
    xrightmargin=4pt
}

% --- Pie de Página Minimalista ---
\usepackage{fancyhdr}
\pagestyle{fancy}
\fancyhf{}
\renewcommand{\headrulewidth}{0pt}
\fancyfoot[C]{\fontsize{7.5pt}{9pt}\selectfont\color{black!70} Pág. \thepage\ de \pageref{LastPage}}

% --- Macro de Diapositiva / Filmina ---
\newcommand{\filmina}[2]{%
  \par\vspace{3.5pt}%
  \noindent\textbf{\sffamily\fontsize{8.5pt}{10.5pt}\selectfont $\blacksquare$\ Diapo #1: #2}%
  \par\vspace{1.5pt}%
}

% --- Macro de Énfasis / Definición Inline ---
\newcommand{\destacado}[1]{%
  \par\vspace{2.5pt}%
  \noindent\textit{\textbf{#1}}%
  \par\vspace{1pt}%
}

\hypersetup{
    colorlinks=false,
    pdfborder={0 0 0}
}

% =========================================================================
% CUERPO DEL DOCUMENTO
% =========================================================================
\begin{document}

% --- Cabecera de 1 sola línea (Cero desperdicio vertical) ---
\noindent
{\fontsize{11pt}{13pt}\selectfont\sffamily\bfseries [ASIGNATURA] -- [UNIDAD / TEMA]}\hfill
{\fontsize{7.5pt}{9.5pt}\selectfont\sffamily\color{black!75} Transcripción de Presentación / Cátedra}\\[1pt]
\rule{\linewidth}{0.6pt}\vspace{1.5mm}

% =========================================================================
% EJEMPLO DE TRANSCRIPCIÓN SECUENCIAL DE FILMINAS
% =========================================================================

\filmina{1}{Carátula e Introducción al Módulo}
\begin{itemize}
    \item Presentación de la unidad temática: Fundamentos de Arquitectura de Sistemas Distribuidos.
    \item Objetivos del módulo: Comprender modelos de consistencia, latencia de red y particionamiento.
    \item Bibliografía obligatoria: Tanenbaum, Cap. 4; Coulouris, Cap. 2 y 3.
\end{itemize}

\filmina{2}{¿Qué es un Sistema Distribuido?}
\begin{itemize}
    \item Colección de computadoras independientes que aparecen ante los usuarios como un sistema único y coherente.
    \item \textbf{Características fundamentales:}
    \begin{itemize}
        \item Ausencia de reloj global compartido.
        \item Memoria físicamente distribuida (paso de mensajes).
        \item Concurrencia inherente de procesos y componentes.
        \item Fallas independientes y parciales de nodos.
    \end{itemize}
\end{itemize}

\filmina{3}{Transparencia en Sistemas Distribuidos}
\begin{itemize}
    \item Ocultar al usuario final la naturaleza distribuida de los recursos:
    \item \textbf{Acceso:} Misma sintaxis para recursos locales y remotos.
    \item \textbf{Ubicación:} El usuario desconoce la máquina física donde reside el dato.
    \item \textbf{Migración:} Los recursos se reubican sin cambiar su nombre lógico.
    \item \textbf{Replicación:} Múltiples copias concurrentes invisibles al cliente.
    \item \textbf{Falla:} Tolerancia y recuperación automática sin aviso al usuario.
\end{itemize}

\filmina{4}{Teorema CAP (Brewer, 2000)}
\destacado{Principio Fundamental de Particionamiento:}
En un sistema distribuido asíncrono con tolerancia a particiones de red, es matemáticamente imposible garantizar simultáneamente más de dos de las siguientes tres propiedades:
\begin{enumerate}
    \item \textbf{Consistency (C):} Todos los nodos leen el valor más reciente en todo momento.
    \item \textbf{Availability (A):} Cada petición sin fallo recibe una respuesta no bloqueante.
    \item \textbf{Partition Tolerance (P):} El sistema sigue funcionando pese a pérdida de mensajes entre nodos.
\end{enumerate}

\filmina{5}{Modelos de Consistencia}
\begin{itemize}
    \item \textbf{Consistencia Estricta:} Toda lectura refleja inmediatamente la escritura más reciente en tiempo absoluto (físicamente irrealizable en sistemas distribuidos).
    \item \textbf{Consistencia Secuencial:} Todos los procesos observan las mismas operaciones en el mismo orden lógico (Lamport).
    \item \textbf{Consistencia Causal:} Las operaciones con relación causal se observan en igual orden; las concurrentes pueden diferir.
    \item \textbf{Consistencia Eventual:} Si cesan las escrituras, todas las réplicas convergen al mismo estado final (Dynamo, Cassandra).
\end{itemize}

\filmina{6}{Ecuación de Tiempo de Comunicación}
El retardo total $T_{\text{msg}}$ para enviar un paquete de $S$ bytes sobre un enlace físico con ancho de banda $B$ y latencia de propagación $L$ se calcula:
$$T_{\text{msg}} = \alpha + \frac{S}{B} + L$$
Donde $\alpha$ es la sobrecarga del software en el emisor y receptor (\textit{software overhead}).

\filmina{7}{Resumen y Puntos Clave de Examen}
\begin{itemize}
    \item En redes WAN, la partición de red ($P$) es inevitable; la elección real en ingeniería siempre es $CP$ vs. $AP$.
    \item Los sistemas $CP$ sacrifican disponibilidad para evitar datos obsoletos (bancos, transacciones ACID).
    \item Los sistemas $AP$ sacrifican consistencia para mantener servicio continuo (redes sociales, DNS, caches).
\end{itemize}

\begin{lstlisting}[style=compactcode, language=bash]
# Monitoreo de latencia y paquetes perdidos entre nodos
ping -c 100 -i 0.2 192.168.20.200 | tail -n 2
mtr --report --report-cycles 10 192.168.20.200
\end{lstlisting}

\end{document}
